% !TEX root = 11.tex
\documentclass[a4paper,10pt]{article}

\usepackage{pdfsync}
\usepackage[utf8]{inputenc}
\usepackage{microtype}
\usepackage{mathtools}
\usepackage{amsthm,amsfonts,amsmath,amssymb}
\usepackage{mathabx}
% \usepackage{MnSymbol}
\usepackage{hyperref}
\usepackage{bbold}
\usepackage[textwidth=2cm,textsize=small]{todonotes}
\usepackage{subcaption}
\usepackage{enumerate}
\usepackage{xspace}
\usepackage{stmaryrd}
\usepackage{verbatim}
\usepackage{cleveref}

\hypersetup{
    colorlinks,
    linkcolor={red!50!black},
    citecolor={blue!50!black},
    urlcolor={blue!80!black}
}

\author{}

\newcommand{\ignore}[1]{}

\newcommand{\st}{s.t.}
\newcommand{\ie}{i.e.}
\newcommand{\eg}{e.g.}

\newcommand{\sep}{\;|\;}
\renewcommand{\rule}[2]{\dfrac{#1}{#2}}
\newcommand{\set}[1]{\left\{#1\right\}}
\newcommand{\setof}[2]{\set{#1 \;\middle|\; #2}}
\newcommand{\tuple}[1]{\left( #1 \right)}
\newcommand{\pair}[2]{\tuple {#1, #2}}
\newcommand{\closure}[3]{\tuple {#1, #2, #3}}
\newcommand{\sem}[1]{\llbracket{#1}\rrbracket}
\newcommand{\powerset}[1]{\mathcal P(#1)}
\newcommand{\map}[2]{\mathsf{map}\; #1\; #2}
\newcommand{\filter}[2]{\mathsf{filter}\; #1\; #2}
\newcommand{\floor}[1]{\left\lfloor #1 \right\rfloor}
\newcommand{\PreRel}[2]{\mathsf{Pre}_{#1}(#2)}
\newcommand{\Pre}[1]{\PreRel {} {#1}}
\newcommand{\PreStar}[1]{\mathsf{Pre}^*(#1)}
\newcommand{\Post}[1]{\mathsf{Post}(#1)}
\newcommand{\PostStar}[1]{\mathsf{Post^*}(#1)}
\newcommand{\Conf}{\mathsf{Conf}}
\newcommand{\upwardclosure}[1]{#1\uparrow}
\newcommand{\card}[1]{\left| #1 \right|}
\newcommand{\lang}[1]{L(#1)}
\newcommand{\goesto}[1]{\xrightarrow{#1}}
\newcommand{\e}{\varepsilon}

\newcommand{\N}{\mathbb N}
\newcommand{\Z}{\mathbb Z}
\newcommand{\Q}{\mathbb Q}

% LTL
\newcommand{\true}{\mathbf{true}}
\newcommand{\false}{\mathbf{false}}
\newcommand{\X}{\mathbf X}
\newcommand{\U}{\mathbf U}
\newcommand{\Release}{\mathbf R}
\newcommand{\F}{\mathbf F}
\newcommand{\G}{\mathbf G}

% CTL
\newcommand{\EX}{\mathbf{EX}}
\newcommand{\AX}{\mathbf{AX}}
\newcommand{\EF}{\mathbf{EF}}
\newcommand{\AF}{\mathbf{AF}}
\newcommand{\EG}{\mathbf{EG}}
\newcommand{\AG}{\mathbf{AG}}
\newcommand{\EU}[2]{\mathbf E(#1 \U #2)}
\newcommand{\AU}[2]{\mathbf A(#1 \U #2)}
\newcommand{\ER}[2]{\mathbf E(#1 \Release #2)}
\newcommand{\AR}[2]{\mathbf A(#1 \Release #2)}

\newcommand{\NL}{\textsf{NL}}
\newcommand{\PTIME}{\textsf{PTIME}}
\newcommand{\PSPACE}{\textsf{PSPACE}}

\makeatletter

\def\dasharrowfill@#1#2#3#4{%
        $\m@th
        \thickmuskip0mu
        \medmuskip\thickmuskip
        \thinmuskip\thickmuskip
        \relax
        #4#1\mkern2mu
        \xleaders\hbox{$#4\mkern2mu#2\mkern2mu$}\hfill
        \mkern2mu
        #3$%
}

\def\dashleftarrowfill@{\dasharrowfill@\leftarrow\relbar\relbar}
\def\dashrightarrowfill@{\dasharrowfill@\relbar\relbar\rightarrow}
\def\dashleftrightarrowfill@{\dasharrowfill@\leftarrow\relbar\rightarrow}
\def\dashLeftarrowfill@{\dasharrowfill@\Leftarrow\Relbar\Relbar}
\def\dashRightarrowfill@{\dasharrowfill@\Relbar\Relbar\Rightarrow}
\def\dashLeftrightarrowfill@{\dasharrowfill@\Leftarrow\Relbar\Rightarrow}

\providecommand*\xdashleftarrow[2][]{%
  \ext@arrow 0055{\dashleftarrowfill@}{#1}{#2}}
\providecommand*\xdashrightarrow[2][]{%
  \ext@arrow 0055{\dashrightarrowfill@}{#1}{#2}}
\providecommand*\xdashleftrightarrow[2][]{%
  \ext@arrow 0055{\dashleftrightarrowfill@}{#1}{#2}}
\providecommand*\xdashLeftarrow[2][]{%
  \ext@arrow 0055{\dashLeftarrowfill@}{#1}{#2}}
\providecommand*\xdashRightarrow[2][]{%
  \ext@arrow 0055{\dashRightarrowfill@}{#1}{#2}}
\providecommand*\xdashLeftrightarrow[2][]{%
  \ext@arrow 0055{\dashLeftrightarrowfill@}{#1}{#2}}

\def\slashedarrowfill@#1#2#3#4#5{%
$\m@th\thickmuskip0mu\medmuskip\thickmuskip\thinmuskip\thickmuskip
  \relax#5#1\mkern-7mu%
  \cleaders\hbox{$#5\mkern-2mu#2\mkern-2mu$}\hfill
  \mathclap{#3}\mathclap{#2}%
  \cleaders\hbox{$#5\mkern-2mu#2\mkern-2mu$}\hfill
  \mkern-7mu#4$%
}

\def\rightslashedarrowfill@{%
  \slashedarrowfill@\relbar\relbar{\raisebox{1.2pt}{$\scriptscriptstyle\diagup$}}\rightarrow}
\newcommand\xslashedrightarrow[2][]{%
  \ext@arrow 0055{\rightslashedarrowfill@}{#1}{#2}}
\makeatother

\newtheorem{exercise}{Exercise}

\newenvironment{solution}{\begin{proof}[\protect{\textnormal{\emph{Solution.}}}]}{\end{proof}}

\let\oldqedhere\qedhere

\newif\ifstartedinmathmode
\renewcommand*{\qedhere}{%
  \relax\ifmmode\startedinmathmodetrue\else\startedinmathmodefalse\fi%
  \ifstartedinmathmode\tag*{\protect\oldqedhere}\else\ifinlist{\hfill\oldqedhere}{\oldqedhere}\fi%
}

\crefname{equation}{}{}
\crefname{exercise}{Exercise}{Exercises}


\title{Formal verification - Tutorial 11}
\date{01/06/2026}

\begin{document}

\maketitle

\section*{Markov chains}

\begin{exercise}
    Let $M = \tuple{S, \to}$ be a finite Markov chain
    and consider a target set $T \subseteq S$.
    %
    For every initial state $s \in S$, let
    %
    \begin{align*}
        x_s = \Pr {M, s \models \F T}
    \end{align*}
    %
    be the probability of reaching a state in $T$ when started from $s$.
    %
    \begin{enumerate}
        \item Can we compute the set of initial states
        $s$ such that $x_s = 0$ without solving linear equations?

        \item Can we compute the set of initial states
        $s$ such that $x_s = 1$ without solving linear equations?
    \end{enumerate}
    %
    Does the answers change if $M$ is infinite-state?
\end{exercise}
\begin{solution}
    Regarding the first question, we note that $x_s = 0$ if, and only if, there is no path from $s$ to any state in $T$.
    Since paths are finite, this observation holds even for infinite-state Markov chains.
    It follows that we can compute the set of states $s$ such that $x_s = 0$
    by performing a graph search on the underlying graph of $M$.

    Regarding the second question, we note that the set of states that appear infinitely often in a run of $M$ is an end component of $M$.
    We can thus decompose $M$ into its end components $C_1, \dots, C_n$ (which are the bottom strongly connected components of the underlying graph).
    We then have that $x_s = 1$ if, and only if, every end component reachable from $s$ intersects $T$.

    The answer for the second question changes for infinite-state Markov chains.
    For instance consider the one-sided random walk
    with reachability probability satisfying $x = 1 - p + p \cdot x^2$.
    %
    Whether $x = 1$ depends on the actual value of $p$, not just on whether $p$ is zero or not.
\end{solution}

\begin{exercise}
    Let $M = \tuple{S, \to}$ be a finite Markov chain
    and consider two sets $I, T \subseteq S$.
    %
    For every initial state $s \in S$, let
    %
    \begin{align*}
        x_s = \Pr {M, s \models I \U T}
    \end{align*}
    %
    be the probability of reaching a state in $T$ while staying in $I$ when started from $s$.
    %
    \begin{enumerate}
        \item Provide a characterisation of the vector of probabilities $\tuple{x_s}_{s \in S}$.
    
        \item Can we compute these probabilities?
    \end{enumerate}
\end{exercise}
\begin{solution}
    We reduce the problem to that of computing reachability probabilities when $I = S$ (trivial invariant).
    First, make all states in $T$ and in $S \setminus (I \cup T)$ absorbing.
    Clearly, this operation does not change the probabilities $x_s$.
    In the new chain, compute the probability of reaching $T$.
    This will be the same as the probability of reaching $T$ while staying in $I$ in the original chain.
    The reason is that in the original chain we can never step out of $I$ while trying to reach $T$.
    In the new chain, we model this by making $T$ unreachable when stepping out of $I$.
\end{solution}

\begin{exercise}
    Let $M$ be a Markov chain
    and let $L$ be an $\omega$-regular property of $M$ recognised by a nondeterministic Büchi automaton $A$.
    %
    Can we decide efficiently
    %
    \begin{enumerate}
        \item whether $\Pr {M \models L} = 0$?
        \item whether $\Pr {M \models L} = 1$?
    \end{enumerate} 
\end{exercise}
\begin{solution}
    First of all, notice that both queries are decidable, in \PSPACE,
    by converting $A$ into a deterministic parity automaton $B$ (with a singly exponential blow-up)
    and then computing the probability that $M$ satisfies the parity condition of $D$ (in \NC 2).

    We show that $\Pr {M \models L} = 1$ is \PSPACE-hard.
    We reduce from the universality problem for nondeterministic finite automata, which is well-known to be \PSPACE-hard.
    Let $B$ be a nondeterministic finite automaton over alphabet $\Sigma = \set {a, b}$.
    We construct a Markov chain $M$ with three states $a, b, \dollar$ and transitions
    %
    \begin{align*}
        &a \goesto {1/3} a, a \goesto {1/3} b, a \goesto {1/3} \dollar, \\
        &b \goesto {1/3} a, b \goesto {1/3} b, b \goesto {1/3} \dollar, \\
        &\dollar \goesto 1 \dollar.
    \end{align*}
    %
    We construct a nondeterministic Büchi automaton $A$
    recognising {$\lang A = \lang B \cdot \dollar^\omega$}.
    %
    We claim that $\Pr {M \models L} = 1$ if, and only if, $B$ is universal ($\lang B = \Sigma^*$).
    %
    The ``only if'' direction is easy: If $B$ is universal, then $\Pr {M \models L} = 1$
    since all runs of $M$ will eventually almost surely reach $\dollar$ and thus satisfy $L$.
    %
    For the ``if'' direction, assume that $B$ is not universal, and let $w \in \Sigma^* \setminus \lang B$.
    We have that the whole cylinder set $w \Sigma^\omega$ is disjoint from $\lang A$.
    This set has measure $1/3^{|w|}$, since $M$ generates words uniformly at random,
    and thus $\Pr {M \models L} < 1 - 1/3^{|w|} < 1$, as required.
\end{solution}

\begin{exercise}
    Let $M$ be a Markov chain
    and let $L$ be an $\omega$-regular property of $M$ recognised by an unambiguous Büchi automaton $A$.
    %
    Show that $\Pr {M \models L}$ can be computed in \PTIME.
\end{exercise}
\begin{solution}
    We pretend that $A$ is deterministic and we compute the product $M \times A$ of $M$ and $A$.
    The product contains a weighted edge
    %
    \begin{align*}
        \tuple {s, p} \goesto x \tuple {t, q}
    \end{align*}
    %
    when $s \goesto x t$ is a transition of $M$ (with probability $x$)
    and $p \goesto s q$ is a transition of $A$ reading input symbol $s$.
    %
    Since $A$ is nondeterministic, the weight of edges outgoing from $\tuple {s, p}$ may sum up to more than $1$.
    Never mind that.
    Decompose $M \times A$ into its bottom strongly connected components
    and let $C$ be the union of all the bottom strongly connected components that contain an accepting state of $A$.
    %
    Consider the least nonnegative solution to the system of equations
    %
    \begin{align*}
        x_{s, p} &= 1 && \text{if $\tuple {s, p} \in C$}, \text{ and otherwise } \\
        x_{s, p} &= \sum_{\tuple {s, p} \goesto x \tuple {t, q}} x \cdot x_{t, q}.
    \end{align*}
    %
    Since $A$ is unambiguous (accepted words have a unique accepting run), $x_{s, p} \in [0, 1]$, and in $\Q$.
    Then, $\Pr {M \models L} = \sum_{s \in S} x_{s, p_0}$, where $p_0$ is the initial state of $A$.
\end{solution}

\section*{Random walks}

\begin{exercise}
    Let $p \in [0, 1]$ be a parameter.
    Consider the random walk induced by the 1-state pPDS
    over stack alphabet $\Gamma = \set X$ and transitions
    %
    \begin{align*}
        X \goesto {1-p} \varepsilon
            \quad\text{and}\quad
                X \goesto p X^3.
    \end{align*}
    %
    Compute the probability of termination from configuration $X$ as a function of $p$.
    For which values of $p$ is this probability equal to $1$?
\end{exercise}
\begin{solution}
    The required probability $x$ is the least nonnegative solution of
    %
    \begin{align*}
        x = 1 - p + p \cdot x^3.
    \end{align*}
    %
    We can solve this equation and obtain that $x$ is the minimum of $1$ and (for $p > 0$)
    %
    \begin{align*}
        \frac { -p + \sqrt{4 p - 3 p^2} } {2 p}
    \end{align*}
    %
    This is a monotonically decreasing function of $p$
    that is equal to $1$ when $p = 1/3$.
\end{solution}

\begin{exercise}
    Let $n \in \N$ and $p \in [0, 1]$ be parameters.
    Consider the random walk induced by the 1-state pPDS
    over stack alphabet $\Gamma = \set X$ and transitions
    %
    \begin{align*}
        X \goesto {1-p} \varepsilon
            \quad\text{and}\quad
                X \goesto p X^n.
    \end{align*}
    %
    Show that the probability of termination $x$ from configuration $X\varepsilon$ is $1$
    if, and only if, $p \leq 1 / n$.
\end{exercise}
\begin{solution}
    The expected drift of the random walk is
    \begin{align*}
        (1-p) \cdot (-1) + p \cdot (n-1) = p \cdot n - 1.
    \end{align*}
    %
    If $p < 1/n$, then the expected drift is negative, and thus $x = 1$.
    If $p > 1/n$, then the expected drift is positive, and thus $x < 1$.
    If $p = 1/n$ (critical value), then the expected drift is zero, and we need to argue differently.
    %
    The required probability $x$ is the least nonnegative solution of
    %
    \begin{align*}
        x = 1 - p + p \cdot x^n.
    \end{align*}
    %
    In other words we have $x^n - \frac 1 p \cdot x + \frac {1 - p} p = 0$.
    The value $x = 1$ is always a solution of this equation, regardless of the value of $p$.
    We thus divide the polynomial by $x - 1$ and obtain
    %
    \begin{align*}
        1 - p \cdot (1 + x + \cdots + x^{n-1}) = 0.
    \end{align*}
    %
    This equation has solution $x = 1$ if, and only if, $p = 1/n$.
\end{solution}

\section*{Probabilistic pushdown automata}

\begin{exercise}
    Show that the quantitative reachability problem for pPDA is decidable with \PSPACE~complexity.
    \emph{Hint: Use the fact that the existential fragment of Tarki's arithmetic is decidable in \PSPACE.}
\end{exercise}
\begin{solution}
    % It suffices to consider separately the three cases $\sim \in \set{<, =, >}$.
    We can massage the logical sentence to have only universal quantification,
    and then we can negate it to obtain an existential sentence.
\end{solution}

\begin{exercise}
    Can we decide efficiently whether $[pXq] = 0$?
    And whether $[pXq] = 1$?
\end{exercise}
\begin{solution}
    The first query $[pXq] = 0$ is decidable in \PTIME.
    Its complement, $[pXq] > 0$, is equivalent to reachability $pX \to^* q\varepsilon$ in the underlying pushdown system,
    which is decidable in \PTIME.
    % TODO: second query $[pXq] = 1$
\end{solution}

\bibliographystyle{plainurl}
\bibliography{bibliography}

\end{document}